\subsection{RELATIONS BETWEEN DERIVATIONS AND ALGEBRA HOMOMORPHISMS}

We suppose again in this no. that Case (II) of no. 2 holds and the graded K-algebra A is not assumed to be associative. Given an element \(\delta \in \Delta\) , consider the graded K-module E( \(\delta\) ) (II, §11, no. 2) such that

\[
(\mathbf {E} (\delta)) _ {\mu} = \mathbf {E} _ {\mu + \delta}
\]

for all \(\mu \in \Delta\) . We define on the graded K-module \(\mathbf{A} \oplus \mathbf{E}(\delta)\) a graded K-algebra structure by setting, for every homogeneous element \(a \in \mathbf{A}\) and arbitrary elements \(a' \in \mathbf{A}, x, x'\) in \(\mathbf{E}(\delta)\)

\[
(a, x) \left(a ^ {\prime}, x ^ {\prime}\right) = \left(a a ^ {\prime}, x. a ^ {\prime} + \varepsilon_ {\delta, \deg (a)} a. x ^ {\prime}\right);\tag{28}
\]

the verification of the fact that this multiplication defines a graded ring structure is immediate.

The projection \(p:(a,x)\mapsto a\) is called the augmentation of the algebra \(\mathbf{A}\oplus\mathbf{E}(\delta)\) and is a graded algebra homomorphism. The graded K-linear mappings \(g:\mathbf{A}\to\mathbf{A}\oplus\mathbf{E}(\delta)\) of degree 0 such that the composition

\[
\mathrm{A} \xrightarrow {g} \mathrm{A} \oplus \mathrm{E} (\delta) \xrightarrow {p} \mathrm{A}
\]

is the identity \(l_{A}\) are the mappings of the form \(x \mapsto (x, f(x))\) , where \(f: A \to E\) is a graded K-linear mapping of degree \(\delta\) .

PROPOSITION 12. For a graded K-linear mapping \(f:A\to E\) of degree \(\delta\) to be an \(\varepsilon\) -derivation, it is necessary and sufficient that the mapping \(x\mapsto(x,f(x))\) of A into \(\mathbf{A}\oplus\mathbf{E}(\delta)\) be a graded K-algebra homomorphism.

Using the fact that for \(x\) homogeneous in \(\mathbf{A}\) and \(y \in \mathbf{A}\)

\[
(x y, f (x y)) = (x, f (x)). (y, f (y)),
\]

we obtain, taking account of (28), the equivalent relation

\[
f (x y) = f (x). y + \varepsilon_ {\delta , \deg (x)} x. f (y),
\]

whence the proposition.

PROPOSITION 13. For the algebra \(\mathbf{A} \oplus \mathbf{E}(\delta)\) to be associative and unital, it is necessary and sufficient that \(\mathbf{A}\) be associative and unital, and that the mappings \((a, x) \mapsto a.x\) and \((a, x) \mapsto x\) . \(a\) define on \(\mathbf{E}\) an \((\mathbf{A}, \mathbf{A})\) -bimodule structure; the unit element of \(\mathbf{A} \oplus \mathbf{E}(\delta)\) is then \((1, 0)\) .

If an element \((u,m)\in\mathbf{A}\oplus\mathbf{E}(\delta)\) is written as unit element of this algebra, it is immediately found that u must be the unit element of A; writing \((u,m).(0,x)=(0,x).(u,m)=(0,x)\) , we obtain u.x=x.u=x for \(x\in E\) and, writing \((u,m).(u,0)=(u,0).(u,m)=(u,0)\) , we obtain m=0. The fact that A is associative when \(\mathbf{A}\oplus\mathbf{E}(\delta)\) is follows from the fact that the augmentation is a surjective homomorphism. The condition \((x.a') .a''=x.(a'.a'')\) is then equivalent to \(((0,x)(a',0))(a'',0)=(0,x)((a',0)(a'',0))\) and similarly the condition \(a.(a'.x)=(a.a').x\) is equivalent to

\[
(a, 0) ((a ^ {\prime}, 0) (0, x)) = ((a, 0) (a ^ {\prime}, 0)) (0, x);
\]

finally the condition \(a \cdot (x \cdot a') = (a \cdot x) \cdot a'\) is equivalent to

\[
(a, 0) ((0, x) (a ^ {\prime}, 0)) = ((a, 0) (0, x)) (a ^ {\prime}, 0).
\]

